\documentclass[11pt,a4paper]{article}
% Préambule commun aux documents de PythonIPM/documentation.
% Compilation : tectonic <fichier>.tex
% tectonic compile avec XeTeX, qui lit l'UTF-8 nativement : surtout PAS inputenc ni fontenc,
% faits pour pdfLaTeX. Ce sont eux qui décodaient mal les caractères multi-octets — « cœur »
% ressortait en « cur », et les tirets cadratins en caractères parasites.
\usepackage{lmodern}
\usepackage[french]{babel}
\usepackage{microtype}
\usepackage[a4paper,margin=2.4cm]{geometry}
\usepackage{amsmath,amssymb}
\usepackage{booktabs}
\usepackage{longtable}
\usepackage{array}
\usepackage{ragged2e}
\usepackage{xcolor}
\usepackage{tcolorbox}
\usepackage{enumitem}
\usepackage{fancyhdr}
\usepackage[hidelinks]{hyperref}

\definecolor{encadre}{HTML}{F4F6F8}
\definecolor{bordure}{HTML}{4A6FA5}
\definecolor{alerte}{HTML}{B03A2E}

% Encadré « à retenir » : la lecture non technique du passage qui précède
\newtcolorbox{retenir}[1][]{
  colback=encadre, colframe=bordure, boxrule=0.6pt, arc=2pt,
  left=8pt, right=8pt, top=6pt, bottom=6pt,
  fonttitle=\bfseries, title={#1}
}

% Encadré d'avertissement : un point ouvert, un écart assumé
\newtcolorbox{attention}[1][]{
  colback=white, colframe=alerte, boxrule=0.6pt, arc=2pt,
  left=8pt, right=8pt, top=6pt, bottom=6pt,
  % le bandeau de titre a déjà la couleur d'alerte en fond : un titre `coltitle=alerte`
  % écrirait du rouge sur du rouge. On garde le blanc par défaut, comme l'encadré « retenir ».
  fonttitle=\bfseries, title={#1}
}

\newcommand{\col}[1]{\texttt{#1}}
% \ensuremath et non \textsubscript : \ipm est employé aussi bien dans le texte que dans les
% formules, et « M\textsubscript{0} » placé en mode mathématique est une erreur — silencieuse en
% PDF, mais qui fait échouer la conversion vers Word.
\newcommand{\ipm}{\ensuremath{M_{0}}}

\setlist[itemize]{itemsep=2pt, topsep=4pt}
\setlist[description]{style=nextline, leftmargin=1.2cm, itemsep=5pt}

% La source du document apparaît en tête de chaque page. Les documents RGPH la redéfinissent
% par \renewcommand{\source}{...} AVANT \begin{document}.
\newcommand{\source}{EHCVM 2021}

\pagestyle{fancy}
\fancyhf{}
\fancyhead[L]{\small\itshape IPM Côte d'Ivoire --- \source}
\fancyhead[R]{\small\thepage}
\renewcommand{\headrulewidth}{0.3pt}

\renewcommand{\arraystretch}{1.15}


\renewcommand{\source}{RGPH 2021 --- variantes Emploi}
\setcounter{tocdepth}{2}

\title{\bfseries Comprendre le calcul de l'IPM pas à pas sur le recensement\\[5pt]
\large Bases, colonnes, code du pipeline, et les trois variantes Emploi\\
SU3, agriculture de subsistance et NEET approché}
\author{Pipeline \col{PythonIPM} --- IPM Côte d'Ivoire}
\date{21 septembre 2026}

\begin{document}
\maketitle

\begin{abstract}
\noindent Cette note est le pendant, pour le \textbf{RGPH 2021}, de la note détaillée déjà
produite sur l'EHCVM. Elle explique, sans supposer de connaissance préalable, comment les
réponses du recensement deviennent des indicateurs de privation au niveau du ménage, quelle
colonne de quelle base sert à chaque indicateur, ce que fait exactement chaque script du
pipeline, puis comment sont calculés l'incidence $H$, l'intensité $A$ et l'indice $M_0$. Elle
détaille enfin les trois variantes demandées pour la dimension Emploi : \textbf{SU3 +
agriculture de subsistance}, \textbf{SU3 uniquement} et \textbf{SU3 + NEET approché}. Les
chiffres portent sur 5\,616\,487 ménages et 27\,984\,405 individus recensés, représentant
29,28 millions de personnes après application du taux net de couverture post-censitaire.
\end{abstract}

\tableofcontents
\newpage

\section{Ce qu'il faut retenir avant d'entrer dans le détail}

Le calcul est strictement le même que sur l'enquête --- c'est la méthode Alkire--Foster --- et
il procède en quatre opérations :

\begin{enumerate}
  \item les réponses individuelles et celles du logement sont ramenées au niveau du ménage ;
  \item chaque situation devient un indicateur binaire : 1 si le ménage est privé, 0 sinon ;
  \item les indicateurs sont pondérés, puis additionnés pour obtenir un score entre 0 et 1 ;
  \item le ménage est pauvre si son score atteint au moins un tiers, soit $k=1/3$.
\end{enumerate}

\[
  \boxed{M_0 = H \times A}
\]

Ce qui change par rapport à l'EHCVM, ce n'est donc pas la méthode, c'est \textbf{ce que la
source sait mesurer}. Le recensement ne pose aucune question de santé, ne demande ni les années
d'études ni le nombre de pièces, mais il recense les décès et il porte un \textbf{statut
d'activité déjà construit par l'INS} qui contient la main-d'{\oe}uvre potentielle. C'est ce
dernier point qui rend les variantes Emploi possibles sur le recensement.

\begin{retenir}[Résultat central]
Sur le RGPH, la variante \textbf{SU3 + agriculture de subsistance} donne $H=43{,}1~\%$,
$A=0{,}4464$ et $M_0=0{,}1924$, contre 0,1664 pour l'IPM national de référence. C'est la plus
défendable des trois. \textbf{SU3 uniquement} fait tomber $M_0$ à 0,1380 et \textbf{SU3 + NEET
approché} à 0,1522, mais dans les deux cas parce que l'on \emph{retire} la privation agricole,
qui est de loin la plus contributive du recensement, et non parce que SU3 serait moins sévère.
\end{retenir}

\begin{attention}[Ne jamais comparer l'IPM RGPH à l'IPM EHCVM]
Les deux indices n'ont pas les mêmes indicateurs : l'IPM national du recensement vaut 0,1664,
celui de l'enquête 0,2863. L'écart ne mesure rien de réel. La seule comparaison légitime passe
par la variante \emph{harmonisée}, bâtie sur le socle d'indicateurs commun aux deux sources.
\end{attention}

\section{Données, unités et pondération de la population}

\subsection{Les trois bases utilisées}

L'extrait RGPH 2021 mis à disposition est réparti sur trois fichiers Stata. Le point à
connaître avant de lire le code est que les variables de logement et de biens figurent dans
\emph{deux} fichiers, mais ne sont complètes que dans un seul.

\begin{center}
\small
\begin{tabular}{@{}p{4.6cm}rp{7.4cm}@{}}
\toprule
\textbf{Fichier} & \textbf{Lignes} & \textbf{Rôle dans l'IPM} \\
\midrule
\col{IPM\_Data\_191125.dta} & 27\,984\,405 individus & démographie, éducation, emploi, état
civil, géographie, pondération, taille du ménage. Il porte aussi le logement et les biens, mais
renseignés pour un tiers des ménages seulement : \textbf{ils ne sont pas lus ici} \\
\col{IPM\_Data\_110924\_men.dta} & 5\,528\,535 ménages & logement, eau, éclairage, cuisson,
sanitaires, biens d'équipement --- complets. Une ligne par chef de ménage \\
\col{MORTALITE\_RP2021\_}\allowbreak\col{TRAITEMENT.dta} & 269\,901 décès & décès survenus dans
les 12 mois précédant le recensement, dont l'âge au décès \\
\bottomrule
\end{tabular}
\end{center}

Après agrégation, la table de travail compte \textbf{5\,616\,487 ménages}.

\subsection{Trois clés de jointure différentes}

C'est une particularité du RGPH, et une source d'erreur si on l'ignore :

\begin{description}
  \item[Individus $\rightarrow$ ménage.] La colonne \col{ID\_Menage} (renommée
  \col{id\_menage}) identifie le ménage. C'est la clé de l'IPM RGPH.

  \item[Ménage $\rightarrow$ fichier logement.] Le fichier ménages n'est pas indexé par
  \col{ID\_Menage} mais par \col{INDIV\_ID}, l'identifiant \emph{de l'individu chef de ménage}.
  La jointure passe donc par la ligne du chef repérée dans le fichier individus
  (\col{P16}~=~1). 87\,952 ménages n'ont pas de ligne correspondante : ils sont comptés non
  privés par convention sur les indicateurs de conditions de vie.

  \item[Ménage $\rightarrow$ fichier décès.] Le fichier des décès ne porte \emph{ni}
  \col{INDIV\_ID} \emph{ni} \col{ID\_Menage}. Le seul lien possible est la clé géographique
  complète en douze champs : \col{REGION}, \col{DEPART}, \col{SOUSPREFID}, \col{P\_ZC},
  \col{P04}, \col{P05}, \col{P06}, \col{P07}, \col{P09}, \col{P09A}, \col{P09B}, \col{P10}.
\end{description}

\subsection{Ménage de construction, individu de comptage}

La privation est construite au niveau du \textbf{ménage}, mais les résultats nationaux portent
sur les \textbf{personnes}. Le poids de population du ménage $i$ est :

\[
  n_i = \underbrace{\col{EW}_i}_{\text{taux net de couverture post-censitaire}}
        \times\ \col{TAILLE\_MENAGE}_i .
\]

Le recensement dénombre l'ensemble de la population ; \col{EW} ne corrige pas un plan de
sondage mais la couverture incomplète du dénombrement, mesurée par l'enquête
post-censitaire. La somme $\sum_i n_i$ vaut 29,28 millions de personnes.

\begin{attention}[Ne pas confondre deux pourcentages]
L'agriculture de subsistance concerne 39,8~\% des \emph{ménages} recensés, mais 46,0~\% de la
\emph{population} pondérée : les ménages agricoles sont plus grands que la moyenne. Les indices
$H$, $A$ et $M_0$ utilisent toujours la pondération population.
\end{attention}

\subsection{La grappe, pour les intervalles de confiance}

Le recensement n'a pas de grappe d'échantillonnage. Pour calculer des erreurs-types, le pipeline
reconstruit une unité primaire à partir de la \textbf{zone de dénombrement}, l'unité de terrain
du recensement :

\[
  \col{grappe} = \col{SOUSPREFID}\times 10^{7} + \col{P\_ZC}\times 10^{4} + \col{P05}.
\]

On obtient ainsi 28\,329 zones, qui servent de strates au calcul de la variance.

\section{Des réponses brutes aux indicateurs de privation}

Pour chaque indicateur $j$, le pipeline construit :

\[
g^0_{ij} =
\begin{cases}
1 & \text{si le ménage } i \text{ est privé selon le seuil } z_j,\\
0 & \text{sinon.}
\end{cases}
\]

Un ménage sans personne appartenant à la tranche d'âge d'un indicateur est classé non privé sur
cet indicateur. Ainsi, les 12,6~\% de ménages sans membre de 17 à 40 ans ne sont jamais privés
au titre de SU3, et les 45,7~\% de ménages sans jeune de 15 à 24 ans ne le sont jamais au titre
du NEET.

\subsection{Les quatre indicateurs d'éducation}

\begin{description}
  \item[Fréquentation scolaire.] Colonne \col{P30A}. Privé si au moins un enfant de 6 à 16 ans
  ne fréquente pas actuellement l'école, c'est-à-dire si \col{P30A}~$\neq$~1 (les modalités 2
  « a déjà fréquenté », 3 « jamais » et 8 « ne sait pas » sont toutes des privations).
  Privation : 21,7~\% des ménages, \textbf{35,6~\% de la population}.

  \item[Année de scolarité.] Colonne \col{P32}, diplôme le plus élevé. \textbf{Écart assumé :}
  le recensement ne demande ni le nombre d'années d'études ni la classe atteinte. Les dix années
  accomplies de l'énoncé national sont approchées par « au moins le BEPC », soit les codes 4 à
  22. Le CEPE (code 2, six années) et le CQP (code 3) restent donc sous le seuil. Privé si aucun
  membre de 17 à 95 ans n'atteint ce niveau. Privation : 65,5~\% des ménages, 63,2~\% de la
  population.

  \item[Alphabétisation.] Colonnes \col{P29A} (sait lire et écrire une langue) et \col{P29\_BA}
  (le français). Est alphabétisé celui pour qui \col{P29A}~=~1 \emph{et} \col{P29\_BA}~=~1.
  Privé si au moins un membre de 17 à 49 ans ne l'est pas. Privation : 58,9~\% des ménages,
  \textbf{70,4~\% de la population}.

  \item[État civil.] Colonne \col{P20}. Seule la modalité 1 « déclaré avec extrait d'acte » est
  non privée : la question de l'EHCVM porte sur la \emph{possession} de l'acte, donc « déclaré
  sans acte » (2) est une privation. Privé si au moins un enfant de 5 à 15 ans est dans ce cas.
  Privation : 9,8~\% des ménages, 15,2~\% de la population.
\end{description}

\begin{attention}[Deux écarts au dofile officiel du RGPH 2021, assumés et documentés]
Sur l'alphabétisation, le dofile officiel lit « aucun membre alphabétisé » (33,2~\% des
ménages) là où l'énoncé national dit « un membre ne sait pas lire ou écrire » (70,4~\%). C'est
l'énoncé littéral qui est retenu, comme dans l'EHCVM, pour que les deux sources ne divergent
pas. Sur l'état civil, le dofile officiel ne compte privés que les codes 3 et 8 (9,7~\% des
ménages) : le pipeline y ajoute le code 2, soit 15,2~\%.
\end{attention}

\subsection{L'unique indicateur de santé : la mortalité}

\begin{description}
  \item[Mortalité dans le ménage.] Fichier des décès, colonne \col{M61A2\_AGE} (âge au décès,
  998 = ne sait pas). Privé si au moins un enfant de moins de 18 ans est décédé dans le ménage.
  Privation : 0,7~\% des ménages, 1,1~\% de la population.
\end{description}

Cet indicateur remplace à lui seul les trois indicateurs de santé de l'EHCVM --- assurance
maladie, insécurité alimentaire et renoncement aux soins --- qu'aucune question du recensement
ne permet de reconstruire. Il porte donc \textbf{à lui seul le poids de 0,25} de la dimension.

\begin{attention}[La fenêtre de mortalité n'est pas celle de l'IPM mondial]
Le RGPH recense les décès des \textbf{12 derniers mois} ; l'IPM mondial retient une fenêtre de
\textbf{5 ans}. La privation est donc mécaniquement beaucoup plus rare ici, et les deux ne se
comparent pas. Conséquence directe sur l'indice : avec un poids nominal de 25~\% mais un taux
de privation de 1,1~\%, la dimension Santé ne contribue qu'à 1,6~\% de $M_0$. Le poids de 0,25
est donc, en pratique, largement inemployé --- et c'est ce qui rend l'IPM RGPH si sensible au
choix des indicateurs d'emploi.
\end{attention}

\subsection{Les six indicateurs de conditions de vie}

Tous lus dans le fichier ménages. Les matériaux adéquats sont définis \emph{par exclusion},
comme dans le dofile officiel du RGPH : tout ce qui n'est pas dans la liste courte est
rudimentaire.

\begin{description}
  \item[Électricité.] Colonne \col{P51}. Privé si l'éclairage principal n'est ni le réseau, ni
  un groupe électrogène, ni le solaire. Population privée : 11,5~\%.

  \item[Logement.] Colonnes \col{P45} (murs), \col{P46} (toit), \col{P47} (sol). Privé si l'un
  au moins des trois est rudimentaire. Sont adéquats : sol ciment, carreau ou moquette ; toit
  tôle, béton ou tuile ; murs semi-dur, géobéton ou dur. Le sol en bois et le mur en tôle sont
  comptés rudimentaires, conformément à la classification DHS. Population privée : 25,0~\%.

  \item[Eau potable.] Colonne \col{P49}. Privé si la source n'est pas améliorée au sens des ODD.
  \textbf{Écart assumé :} le temps de trajet jusqu'au point d'eau n'est pas demandé au
  recensement, seul le premier critère du PNUD s'applique --- l'indicateur est donc moins sévère
  que son homologue EHCVM. Population privée : 14,3~\%.

  \item[Énergie de cuisson.] Colonne \col{P52}. Privé si le combustible n'est ni le gaz ni
  l'électricité. Population privée : 62,7~\%.

  \item[Toilettes.] Colonne \col{P48}. Privé si les sanitaires ne sont pas améliorés.
  \textbf{Écart assumé :} le partage des sanitaires avec d'autres ménages n'est pas demandé.
  Population privée : 45,1~\%.

  \item[Biens d'équipement.] Colonnes \col{P55AA} à \col{P57EE}. Privé si le ménage possède au
  plus un bien parmi radio, télévision, téléphone (fixe ou mobile), ordinateur, charrette, vélo,
  moto et réfrigérateur, et ne possède pas de voiture (\col{P55CC}). \textbf{Ici le recensement
  est plus complet que l'enquête :} les 8 biens de la définition PNUD y sont tous, charrette
  (« véhicule à traction animale ») comprise, alors que l'EHCVM n'en a que 7. Population
  privée : 18,2~\%.
\end{description}

\begin{retenir}[Ce que le recensement ne sait pas mesurer]
Quatre indicateurs de l'EHCVM disparaissent : assurance maladie, insécurité alimentaire,
renoncement aux soins (aucune question de santé) et promiscuité (le nombre de pièces n'est pas
collecté). L'IPM national du RGPH compte donc \textbf{13 indicateurs} contre 16 dans l'enquête,
et sa dimension Santé n'a plus qu'un seul indicateur.
\end{retenir}

\section{Les indicateurs d'emploi, dans le détail}

C'est la section qui justifie cette note : les trois variantes demandées ne pouvaient, jusqu'à
présent, être calculées que sur l'EHCVM. Elles le sont désormais aussi sur le recensement.

\subsection{Pourquoi SU3 est constructible sur le RGPH}

L'EHCVM permet de reconstruire la main-d'{\oe}uvre potentielle à la main, à partir de la
recherche d'emploi, de la disponibilité et du désir de travailler (questions 4.13 à 4.18). Le
recensement ne pose pas ces questions dans le même détail --- et c'est ce qui avait fait écarter
les variantes Emploi côté RGPH.

Mais le recensement porte une variable \textbf{déjà construite par l'INS},
\col{Statut\_OQPtbb}, qui code directement le statut d'activité en quatre modalités :

\begin{center}
\small
\begin{tabular}{@{}clrr@{}}
\toprule
\textbf{Code} & \textbf{Statut} & \textbf{Individus} & \textbf{Part} \\
\midrule
0 & Occupé & 11\,082\,635 & 39,6~\% \\
1 & Chômeur au sens du BIT & 403\,834 & 1,4~\% \\
2 & Inactif & 13\,794\,677 & 49,3~\% \\
\textbf{3} & \textbf{Main-d'{\oe}uvre potentielle} & \textbf{2\,127\,898} & \textbf{7,6~\%} \\
--- & Non renseigné (hors champ d'âge) & 575\,361 & 2,1~\% \\
\midrule
 & Total & 27\,984\,405 & 100,0~\% \\
\bottomrule
\end{tabular}
\end{center}

La modalité 3 \emph{est} la main-d'{\oe}uvre potentielle. SU3 se lit donc directement :

\[
  \text{SU3} = \{\col{Statut\_OQPtbb} \in \{1,\,3\}\}
  = \text{chômage BIT} \ \cup\ \text{main-d'{\oe}uvre potentielle}.
\]

\begin{retenir}[Le point méthodologique de cette note]
SU3 n'est pas \emph{approché} sur le RGPH : il est \textbf{lu} dans une variable officielle de
l'INS, là où l'EHCVM doit le reconstruire. Le recensement est, sur ce point précis, la source la
plus directe des deux.
\end{retenir}

Le choix de \col{Statut\_OQPtbb} plutôt que de \col{Statut\_OQPtb}, l'autre variable de statut
présente dans la base, n'est pas anodin : \col{Statut\_OQPtb} donne 21,6~\% de ménages privés au
chômage contre 5,9~\% ici, parce qu'elle n'intègre pas le reclassement INS des personnes
déclarées sans activité mais en réalité occupées. \col{Statut\_OQPtbb} a par ailleurs été
vérifiée contre la construction du dofile officiel du RGPH 2021, qui applique les trois critères
BIT explicitement : les deux ne diffèrent que sur 18\,910 individus sur 28 millions, soit
0,07~\%.

\subsection{Chômage BIT ou SU1}

Privé si au moins un membre de 17 à 40 ans a \col{Statut\_OQPtbb}~=~1. On compte 298\,644
chômeurs dans cette tranche sur 11,4 millions de personnes concernées, soit \textbf{2,6~\%}. Au
niveau du ménage : 196\,019 ménages privés, 3,5~\% des ménages et \textbf{5,9~\% de la
population}.

Le chômage BIT est rare, pour les mêmes raisons que dans l'enquête : une heure de travail suffit
à être classé occupé, et les personnes découragées qui ne recherchent plus activement sortent du
chômage.

\subsection{SU3 : chômage plus main-d'{\oe}uvre potentielle}

Privé si au moins un membre de 17 à 40 ans a \col{Statut\_OQPtbb}~$\in \{1,3\}$. On compte
1\,661\,331 personnes, soit \textbf{14,6~\%} des 17--40 ans --- cinq fois et demie le chômage
BIT. Au niveau du ménage : 1\,033\,121 ménages, 18,4~\% des ménages et \textbf{26,3~\% de la
population}.

Tout chômeur BIT appartient nécessairement à SU3. La propriété a été vérifiée dans les données :
\textbf{aucun} des 196\,019 ménages privés au titre du chômage ne se trouve hors de SU3.

\begin{attention}[Pourquoi il ne faut pas additionner mécaniquement SU1 et SU3]
SU3 contient déjà SU1. Une variante « SU1 + SU3 » attribue donc deux privations au même chômage.
Ce montage peut être calculé comme test (option \col{-{}-su1\_su3}), mais il n'est pas recommandé
comme spécification principale.
\end{attention}

\subsection{Agriculture de subsistance : ce que mesure réellement le code}

Colonnes \col{Statut\_OQPtbb}, \col{P36BNEWtb} (branche d'activité, nomenclature NAEMA) et
\col{P37NEWtb} (situation dans la profession). Le chef de ménage est compté en
subsistance s'il
est simultanément :

\begin{itemize}
  \item \textbf{occupé} : \col{Statut\_OQPtbb}~=~0 ;
  \item dans une \textbf{branche agricole} : \col{P36BNEWtb} dans l'intervalle
  $[01000\,;\,04000[$, soit les divisions NAEMA 01 à 03 (agriculture, sylviculture, pêche) ;
  \item \textbf{à son compte ou comme aide familial} : \col{P37NEWtb} $\in \{3, 5\}$, ce qui
  exclut le salariat, l'apprentissage et le statut d'employeur.
\end{itemize}

Le ménage n'est privé que si c'est le \textbf{chef} qui est dans ce cas (\col{P16}~=~1). Cette
règle donne 2\,234\,464 ménages privés, soit 39,8~\% des ménages et \textbf{46,0~\% de la
population}.

\begin{attention}[Une approximation différente de celle de l'EHCVM]
L'EHCVM observe le travail \emph{exclusif sur son propre champ} ; le RGPH ne pose pas cette
question et observe la \emph{branche} croisée avec le \emph{statut non salarié}. Les deux
approximations visent la même réalité mais ne coïncident pas : la version RGPH inclut la pêche
et la sylviculture, et n'exclut pas un chef qui ferait aussi du commerce. Dans les deux cas, la
destination de la production --- le vrai critère OIT de la subsistance --- n'est pas observée.
Le libellé prudent reste : « chef de ménage dépendant d'une activité agricole non salariée ---
proxy de l'agriculture de subsistance ».
\end{attention}

\subsection{NEET approché : en réalité un NEE}

L'indicateur international NEET porte sur les jeunes de 15 à 24 ans qui ne sont ni en emploi, ni
en éducation, ni en formation. Le recensement observe les deux premiers termes, pas le
troisième. Le pipeline construit donc, sur les colonnes \col{Statut\_OQPtbb} et \col{P30A} :

\[
  \text{NEE approché} = (15\text{--}24\text{ ans}) \cap
  (\col{Statut\_OQPtbb} \neq 0) \cap (\col{P30A} \neq 1).
\]

On obtient 1\,692\,665 jeunes sur 5\,920\,170, soit \textbf{28,6~\%}. Au niveau du ménage :
1\,105\,964 ménages, 19,7~\% des ménages et \textbf{30,9~\% de la population}. Comme dans
l'EHCVM, le chiffre est une \textbf{borne supérieure} du NEET : des jeunes comptés NEE peuvent
suivre une formation non observée. L'indicateur garde le nom \col{neet\_approx} précisément pour
que cette limite ne se perde pas en route.

SU3 et NEE se recouvrent dans 473\,999 ménages, soit 14,2~\% de la population. Le couple
SU3~+~NEE mesure deux objets différents, mais donne deux fois du poids à certains épisodes de
non-emploi des jeunes.

\section{Que fait chaque code du pipeline}

Le pipeline tourne sur \textbf{deux sources} avec le même enchaînement d'étapes. La source est
portée par la variable d'environnement \col{IPM\_SOURCE} (\col{ehcvm} ou \col{rgph}) ; elle
détermine le dossier de données, la clé du ménage, les désagrégations disponibles et le
\emph{suffixe} de tous les fichiers produits. Les sorties des deux sources cohabitent donc sans
se marcher dessus.

\subsection{Vue d'ensemble}

\begin{center}
\scriptsize
\begin{tabular}{@{}p{4.3cm}p{5cm}p{5.4cm}@{}}
\toprule
\textbf{Script} & \textbf{Entrée} & \textbf{Sortie} \\
\midrule
\col{orchestrateur.py} & --- & socle commun (chemins, clés, journaux) \emph{et} chef d'orchestre \\
\col{01\_rgph\_preconstruction\_}\allowbreak\col{matrice\_situationnelle.py} & les 3 bases
  \col{.dta} du recensement & \col{preconstruction\_matrice\_}\allowbreak\col{situationnelle\_rgph2021.dta} \\
\col{02\_construction\_matrice\_}\allowbreak\col{situationnelle.py} & la préconstruction &
  \col{matrice\_situationnelle\_}\allowbreak\col{rgph2021.dta} + le vecteur $z$ \\
\col{03\_matrice\_privations\_}\allowbreak\col{ponderees.py} & la matrice $g^0$ & poids $w$,
  matrice pondérée, score $c_i$ \\
\col{04\_matrice\_privations\_}\allowbreak\col{censuree.py} & la matrice pondérée & matrice
  censurée, statuts au seuil $k$ \\
\col{05\_indices\_ipm.py} & la matrice $g^0$ + le vecteur $z$ & $H$, $A$, $M_0$, contributions,
  désagrégations, classeurs \col{.xlsx} \\
\col{06\_validation\_afmpi.py} & les sorties de 05 & contrôle croisé par un package indépendant \\
\bottomrule
\end{tabular}
\end{center}

\subsection{\col{orchestrateur.py} --- le socle et le chef d'orchestre}

Ce module a deux rôles. Comme \textbf{socle}, il porte ce qui est commun à toutes les étapes :
les chemins du projet, la clé de fusion (\col{id\_menage} pour le RGPH,
\col{grappe~+~menage~+~vague} pour l'EHCVM), les colonnes techniques transportées d'une étape à
l'autre, la journalisation, et deux listes décisives :

\begin{description}
  \item[\col{SOCLE\_COMMUN}] les indicateurs que l'EHCVM \emph{et} le RGPH savent mesurer. Il
  sert à la variante harmonisée, seule comparable en niveau d'une source à l'autre.

  \item[\col{HORS\_IPM\_NATIONAL}] les indicateurs produits par l'étape 02 mais volontairement
  \emph{exclus} de l'indice publié : \col{chomage\_su3} et \col{neet\_approx}. Ils existent dans
  la matrice de privation et dans le vecteur $z$, et seules les variantes qui les nomment les
  emploient. C'est ce mécanisme qui garantit que l'ajout de SU3 et du NEET au recensement
  \textbf{n'a pas modifié d'un millième} l'IPM national : il vaut 0,1664 avant comme après.
\end{description}

Comme \textbf{chef d'orchestre}, il enchaîne les étapes en sous-processus --- une étape qui
échoue arrête la chaîne --- et écrit un journal d'ensemble. Il porte aussi une fonction
\col{compacter} qui ramène chaque colonne entière au plus petit type qui la contient : sur 5,6
millions de lignes, stocker un indicateur 0/1 en \col{int64} octuple la taille des fichiers de
travail. Les vrais flottants (scores $c_i$, pondérations) restent en \col{float64}, car ils sont
comparés d'une étape à l'autre à $10^{-9}$ près.

\subsection{\col{01\_rgph\_preconstruction\_matrice\_situationnelle.py}}

C'est le seul script réellement spécifique au recensement. Son contrat de sortie : \textbf{une
ligne par ménage, et pour chaque indicateur l'effectif \emph{concerné} à côté de l'effectif
\emph{défavorable}} --- sans qu'aucun seuil ne soit appliqué à ce stade. Il procède en cinq
temps :

\begin{enumerate}
  \item \textbf{Parcours du fichier individus par morceaux de 3 millions de lignes.} 28 millions
  de lignes ne tiennent pas en mémoire. Chaque morceau est agrégé séparément par ménage, puis
  les agrégats sont resommés : un ménage coupé entre deux morceaux est correctement reconstitué,
  puisque la somme d'une somme reste une somme. Les lignes de chef sont mises de côté au passage,
  car elles portent la géographie, la pondération et la taille du ménage.

  \item \textbf{Situations individuelles} (fonction \col{situations\_individuelles}). Chaque
  colonne est un booléen, et l'appartenance à la tranche d'âge concernée est portée par une
  colonne \emph{séparée} --- c'est elle qui permettra à l'étape 02 de distinguer « non privé »
  de « non concerné ». Un âge « ne sait pas » (code 999) ne fait entrer dans aucune tranche.
  \textbf{C'est ici qu'ont été ajoutées} les colonnes \col{chomeurs\_su3\_17\_40},
  \col{jeunes\_15\_24} et \col{jeunes\_neet\_approx\_15\_24}.

  \item \textbf{Conditions de vie et biens} (fichier ménages), avec le comptage des 8 biens
  PNUD.

  \item \textbf{Mortalité juvénile} (fichier des décès), jointe par la clé géographique en douze
  champs.

  \item \textbf{Assemblage et contrôles} : unicité de la clé, pondérations strictement
  positives, distributions, valeurs manquantes.
\end{enumerate}

\subsection{\col{02\_construction\_matrice\_situationnelle.py} --- où sont écrites les règles}

À partir de cette étape, le code est \textbf{commun aux deux sources} : il ne travaille plus que
sur des colonnes de situation normalisées. Ce qui change d'une source à l'autre, ce sont les
énoncés et les seuils, rassemblés dans une carte officielle des indicateurs. Chaque entrée
déclare :

\begin{description}
  \item[\col{colonne}] le nom produit dans $X$ (1 = privé, 0 = non privé) ;
  \item[\col{regle}] le seuil appliqué, écrit comme une fonction ;
  \item[\col{situation}] les colonnes de la préconstruction qui portent la situation brute ;
  \item[\col{eligibilite}] la colonne qui dit combien de membres sont concernés.
\end{description}

Chaque indicateur suit soit la \emph{proposition nationale}, soit l'\emph{application du PNUD},
et l'étiquette est portée dans la carte. Le script produit aussi le \textbf{vecteur $z$}, la
liste des indicateurs avec leur dimension et leur valeur de coupure : c'est lui que lira
l'étape 05 pour construire les poids, ce qui garantit que $w$ et $z$ portent toujours sur
exactement les mêmes colonnes.

\textbf{C'est ici qu'ont été ajoutés} les deux indicateurs \col{chomage\_su3} et
\col{neet\_approx} à la liste du RGPH, qui passe de 13 à 15 indicateurs --- dont 13 seulement
entrent dans l'IPM national.

\subsection{\col{03\_matrice\_privations\_ponderees.py}}

Calcule le vecteur de poids $w$ : les dimensions sont équipondérées, et à l'intérieur d'une
dimension le poids est partagé également entre les indicateurs. Puis la matrice pondérée
$w_j g^0_{ij}$ et le score $c_i$. \textbf{C'est ce script qui retire \col{HORS\_IPM\_NATIONAL}}
du vecteur $z$ avant de calculer les poids : la chaîne 03--04 ne voit donc jamais SU3 ni le NEET,
et l'IPM national reste celui qui est publié.

\subsection{\col{04\_matrice\_privations\_censuree.py}}

Applique le seuil $k=1/3$ et remet à zéro les privations des ménages \emph{non} pauvres. C'est
la censure Alkire--Foster, qui rend $M_0$ décomposable par indicateur. Le script calcule aussi
les deux statuts complémentaires : vulnérable ($0{,}20 \leq c_i < 1/3$) et pauvreté sévère
($c_i \geq 0{,}50$).

\subsection{\col{05\_indices\_ipm.py} --- les variantes}

Le script principal. Il repart de la matrice $g^0$ brute et \textbf{refait} la pondération et la
censure pour chaque variante, puisque les poids changent d'une variante à l'autre. Trois
variantes sont toujours calculées :

\begin{description}
  \item[IPM national] 4 dimensions, tous les indicateurs sauf \col{HORS\_IPM\_NATIONAL} ;
  \item[IPM international] 3 dimensions, sans l'Emploi ;
  \item[IPM harmonisé] le socle commun EHCVM/RGPH, seule variante comparable entre sources.
\end{description}

La variante nationale est comparée par \col{assert} à la matrice censurée de l'étape 04 : si les
deux chemins de calcul divergeaient d'un millième, le script s'arrêterait. C'est ce contrôle qui
garantit qu'aucune modification du code ne déplace silencieusement l'indice publié.

À ces trois variantes s'ajoutent, \textbf{sur option}, les variantes Emploi. Chacune est définie
par la liste des indicateurs qu'elle \emph{exclut} :

\begin{center}
\small
\begin{tabular}{@{}llp{6.4cm}@{}}
\toprule
\textbf{Option} & \textbf{Emploi retenu} & \textbf{Exclusions déclarées} \\
\midrule
\col{-{}-su1\_su3} & chômage, SU3, subsistance & \col{neet\_approx} \\
\col{-{}-su3} & SU3, subsistance & \col{chomage}, \col{neet\_approx} \\
\col{-{}-su3\_seul} & SU3 & \col{chomage}, \col{neet\_approx}, \col{emploi\_subsistance} \\
\col{-{}-su3\_neet} & SU3, NEET approché & \col{chomage}, \col{emploi\_subsistance} \\
\bottomrule
\end{tabular}
\end{center}

Les exclusions sont \emph{explicites} plutôt qu'implicites, afin qu'un indicateur expérimental
ne puisse pas contaminer une variante par oubli. Le script produit ensuite, pour chaque
variante, les indices désagrégés par milieu, région, département, sous-préfecture et taille du
ménage, avec intervalles de confiance à 95~\% et coefficient de variation, puis les
contributions de chaque indicateur à $M_0$, et enfin un classeur Excel de restitution.

\subsection{\col{06\_validation\_afmpi.py}}

Hors chaîne. Recalcule $H$, $A$, $M_0$ et les contributions avec le package indépendant
\col{afmpi}, et compare à $10^{-6}$ près. C'est un contrôle croisé : deux implémentations
différentes de la même méthode doivent donner le même nombre.

\subsection{Comment lancer}

\begin{center}
\small
\begin{tabular}{@{}p{8.6cm}p{6.4cm}@{}}
\toprule
\textbf{Commande} & \textbf{Effet} \\
\midrule
\col{python pipeline/orchestrateur.py} & les deux sources, tout le pipeline \\
\col{python pipeline/orchestrateur.py -{}-source rgph} & le recensement seul \\
\col{python pipeline/orchestrateur.py -{}-check} & les auto-contrôles de chaque étape \\
\col{python pipeline/orchestrateur.py -{}-su3 -{}-su3\_seul -{}-su3\_neet} & ajoute les trois
  variantes Emploi, sur les deux sources \\
\bottomrule
\end{tabular}
\end{center}

\section{Ce qui a été modifié pour rendre les variantes possibles sur le RGPH}

Quatre fichiers, aucune modification de la méthode de calcul :

\begin{description}
  \item[Étape 01 RGPH --- \col{01\_rgph\_preconstruction\_...py}] ajout de la tranche d'âge
  NEET (15--24 ans), de la constante \col{SOUS\_UTILISATION\_SU3}~=~$\{1,3\}$, et des trois
  colonnes de situation \col{chomeurs\_su3\_17\_40}, \col{jeunes\_15\_24} et
  \col{jeunes\_neet\_}\allowbreak\col{approx\_15\_24}. Les auto-contrôles unitaires ont été étendus en
  conséquence : un inactif n'entre ni dans SU1 ni dans SU3, une personne de la main-d'{\oe}uvre
  potentielle entre dans SU3 mais pas dans SU1, un jeune scolarisé n'est pas NEET.

  \item[Étape 02 --- \col{02\_construction\_matrice\_...py}] ajout des deux indicateurs
  \col{chomage\_su3} et \col{neet\_approx} à la carte du RGPH, avec leurs énoncés propres au
  recensement. La source passe de 13 à 15 indicateurs produits.

  \item[\col{orchestrateur.py}] suppression du filtre qui retirait les options Emploi du passage
  sur le recensement. Les quatre options valent désormais pour les deux sources.

  \item[\col{05\_indices\_ipm.py}] suppression du garde-fou \col{assert SOURCE == "ehcvm"} qui
  interdisait les variantes Emploi au RGPH.
\end{description}

\begin{retenir}[L'IPM national n'a pas bougé]
\col{chomage\_su3} et \col{neet\_approx} figurent dans \col{HORS\_IPM\_NATIONAL} : l'étape 03
les retire du vecteur $z$ et la variante nationale de l'étape 05 les exclut explicitement. Après
recalcul complet, l'IPM national du RGPH vaut toujours $H = 38{,}8~\%$, $A = 0{,}4287$ et
$M_0 = 0{,}1664$, identiques aux valeurs publiées dans la note de lecture des sorties.
\end{retenir}

\section{Pondérations : comment chaque privation entre dans le score}

Les quatre dimensions sont équipondérées : chacune reçoit 0,25. À l'intérieur d'une dimension,
ce poids est partagé également entre les indicateurs retenus.

\subsection{IPM national de référence}

\begin{center}
\small
\begin{tabular}{@{}lrrr@{}}
\toprule
\textbf{Dimension} & \textbf{Nombre d'indicateurs} & \textbf{Poids par indicateur} &
\textbf{Poids total} \\
\midrule
Éducation & 4 & 0,0625 & 0,25 \\
Santé & 1 & 0,2500 & 0,25 \\
Emploi & 2 & 0,1250 & 0,25 \\
Conditions de vie & 6 & 0,041667 & 0,25 \\
\midrule
Total & 13 & --- & 1,00 \\
\bottomrule
\end{tabular}
\end{center}

Le contraste avec l'EHCVM est frappant : là où l'enquête partage les 0,25 de la Santé entre
trois indicateurs, le recensement les concentre sur la seule mortalité.

\subsection{Poids dans les trois variantes demandées}

\begin{center}
\small
\begin{tabular}{@{}p{4.3cm}p{7cm}p{2.1cm}@{}}
\toprule
\textbf{Variante} & \textbf{Indicateurs d'emploi retenus} & \textbf{Poids} \\
\midrule
SU3 + agriculture & SU3 ; chef en activité agricole non salariée & 0,125 chacun \\
SU3 uniquement & SU3 & 0,250 \\
SU3 + NEET approché & SU3 ; NEE des 15--24 ans & 0,125 chacun \\
\bottomrule
\end{tabular}
\end{center}

Les poids des trois autres dimensions ne changent jamais. Seul le partage des 0,25 de la
dimension Emploi varie.

\section{Du score du ménage à la pauvreté multidimensionnelle}

\subsection{Score pondéré, seuil, censure}

Pour le ménage $i$ :

\[
  c_i = \sum_{j=1}^{J} w_j g^0_{ij},
  \qquad
  \rho_i(k) = \mathbf{1}(c_i \geq 1/3),
  \qquad
  g^0_{ij}(k) = g^0_{ij}\,\rho_i(k).
\]

Un score de 0,40 signifie que le ménage subit 40~\% de l'ensemble pondéré des privations
possibles. Il ne signifie pas que 40~\% des membres du ménage sont privés. Un score exactement
égal à un tiers est pauvre.

\subsection{H, A et M0}

Soit $n_i$ le poids de population du ménage et $q_i=\rho_i(k)$ son statut de pauvreté :

\[
  H = \frac{\sum_i n_i q_i}{\sum_i n_i},
  \qquad
  A = \frac{\sum_i n_i c_i q_i}{\sum_i n_i q_i},
  \qquad
  M_0 = H A = \frac{\sum_i n_i c_i q_i}{\sum_i n_i}.
\]

\begin{attention}[M0 n'est pas un pourcentage ordinaire]
Un $M_0$ de 0,1924 ne signifie pas que 19,24~\% des personnes sont pauvres. Dans la variante
SU3~+~agriculture, 43,1~\% sont pauvres et leur intensité moyenne est de 44,64~\%. Le produit
$0{,}4311 \times 0{,}4464$ donne 0,1924.
\end{attention}

\section{Exemple entièrement calculé pour un ménage fictif}

Considérons la variante SU3 + agriculture sur le RGPH. Supposons un ménage privé sur six
indicateurs : année de scolarité, alphabétisation, SU3, agriculture, énergie de cuisson et
toilettes --- mais \emph{sans} décès d'enfant.

\begin{center}
\small
\begin{tabular}{@{}lccc@{}}
\toprule
\textbf{Indicateur privé} & $g^0_{ij}$ & $w_j$ & $w_jg^0_{ij}$ \\
\midrule
Année de scolarité & 1 & 0,062500 & 0,062500 \\
Alphabétisation & 1 & 0,062500 & 0,062500 \\
SU3 & 1 & 0,125000 & 0,125000 \\
Agriculture de subsistance & 1 & 0,125000 & 0,125000 \\
Énergie de cuisson & 1 & 0,041667 & 0,041667 \\
Toilettes & 1 & 0,041667 & 0,041667 \\
\midrule
Score $c_i$ & & & \textbf{0,458334} \\
\bottomrule
\end{tabular}
\end{center}

Comme $0{,}458334 \geq 1/3$, le ménage est pauvre. Dans la variante \textbf{SU3 uniquement}, la
privation agricole disparaît et SU3 vaut 0,25 : le score devient
$0{,}0625+0{,}0625+0{,}25+0{,}041667+0{,}041667 = 0{,}458334$, identique par coïncidence. Mais
pour un ménage agricole \emph{sans} SU3 --- le cas le plus fréquent en milieu rural --- le score
perd 0,125 d'un coup et peut repasser sous le seuil. C'est tout le mécanisme de la baisse
observée.

\section{Résultats nationaux des variantes}

\begin{center}
\small
\begin{tabular}{@{}p{4.8cm}rrrrr@{}}
\toprule
\textbf{Dimension Emploi} & $H$ & $A$ & $M_0$ & \textbf{Vulnér.} & \textbf{Sévère} \\
\midrule
Référence : chômage BIT + agriculture & 38,8~\% & 0,4287 & 0,1664 & 20,7~\% & 7,8~\% \\
SU3 + agriculture de subsistance & 43,1~\% & 0,4464 & \textbf{0,1924} & 20,8~\% & 12,3~\% \\
SU3 uniquement & 30,9~\% & 0,4473 & \textbf{0,1380} & 32,6~\% & 9,2~\% \\
SU3 + NEET approché (NEE) & 34,6~\% & 0,4401 & \textbf{0,1522} & 28,3~\% & 9,0~\% \\
\bottomrule
\end{tabular}
\end{center}

Les intervalles de confiance à 95~\% de $M_0$ sont respectivement :

\begin{itemize}
  \item référence : $[0{,}1645\ ;\ 0{,}1684]$ ;
  \item SU3 + agriculture : $[0{,}1903\ ;\ 0{,}1946]$ ;
  \item SU3 uniquement : $[0{,}1366\ ;\ 0{,}1395]$ ;
  \item SU3 + NEE : $[0{,}1506\ ;\ 0{,}1537]$.
\end{itemize}

Les intervalles sont environ dix fois plus étroits que ceux de l'EHCVM : c'est l'effet de
l'exhaustivité du recensement. \textbf{Aucun des écarts entre variantes n'est imputable à
l'imprécision statistique} --- ils sont tous largement significatifs.

\subsection{Pourquoi SU3 + agriculture augmente l'IPM de 16~\%}

SU3 remplace le chômage BIT au même poids de 0,125, mais la privation population passe de
5,9~\% à 26,3~\%, soit un quadruplement. De nombreux ménages auparavant juste sous un tiers
franchissent le seuil : $H$ passe de 38,8~\% à 43,1~\%. Chez les pauvres, SU3 ajoute aussi de la
privation, ce qui porte $A$ de 0,4287 à 0,4464. La pauvreté sévère, elle, augmente de moitié
(7,8~\% $\rightarrow$ 12,3~\%).

\subsection{Pourquoi SU3 uniquement fait fortement baisser l'IPM}

SU3 reçoit 0,25, mais la privation agricole --- qui concernait 46,0~\% de la population, la plus
fréquente des indicateurs d'emploi --- disparaît. La hausse du poids de SU3 ne compense pas
cette suppression, parce que les deux privations ne touchent pas les mêmes ménages : un chef
agriculteur est \emph{occupé}, donc hors SU3 par construction. Des ménages ruraux perdent 0,125
de score et repassent sous le seuil.

Le signe le plus net est le report vers la vulnérabilité : la part de vulnérables bondit de
20,7~\% à 32,6~\%. Ce ne sont pas des ménages qui vont mieux, ce sont les mêmes ménages passés
juste sous la barre.

\begin{attention}[Une baisse de définition, pas une amélioration]
Les ménages et la date sont identiques dans les quatre colonnes. Seule la définition de l'indice
change. Une comparaison entre variantes n'est en aucun cas une évolution temporelle.
\end{attention}

\subsection{Pourquoi SU3 + NEE donne un résultat intermédiaire}

Le NEE est plus fréquent que SU3 : 30,9~\% de la population au niveau ménage contre 26,3~\%. Il
remplace donc une partie du pouvoir discriminant de l'agriculture, mais déplace la dimension
vers la jeunesse et la déconnexion de l'école. La variante devient moins rurale et plus sensible
aux ménages comprenant des jeunes de 15 à 24 ans. Elle ignore en revanche complètement les
45,7~\% de ménages qui n'en comptent aucun.

\section{Lecture par milieu de résidence}

Le RGPH distingue trois milieux, là où l'EHCVM n'en code que deux.

\begin{center}
\small
\begin{tabular}{@{}p{4.2cm}rrrrrr@{}}
\toprule
& \multicolumn{2}{c}{\textbf{Urbain}} & \multicolumn{2}{c}{\textbf{Rural}} &
\multicolumn{2}{c}{\textbf{Semi-urbain}} \\
\cmidrule(lr){2-3}\cmidrule(lr){4-5}\cmidrule(l){6-7}
\textbf{Variante} & $H$ & $M_0$ & $H$ & $M_0$ & $H$ & $M_0$ \\
\midrule
Référence & 5,4~\% & 0,0212 & 71,1~\% & 0,3075 & 41,0~\% & 0,1705 \\
SU3 + agriculture & 8,9~\% & 0,0355 & 75,4~\% & 0,3425 & 48,3~\% & 0,2079 \\
SU3 uniquement & 17,1~\% & 0,0721 & 43,3~\% & 0,1977 & 35,6~\% & 0,1614 \\
SU3 + NEE & 15,1~\% & 0,0623 & 51,9~\% & 0,2323 & 43,1~\% & 0,1893 \\
\bottomrule
\end{tabular}
\end{center}

Part de population : urbain 44,8~\%, rural 45,8~\%, semi-urbain 9,4~\%.

Le tableau est le résultat le plus important de cette note. Retirer l'agriculture
\textbf{triple} l'IPM urbain (0,0212 $\rightarrow$ 0,0721) et \textbf{réduit d'un tiers} l'IPM
rural (0,3075 $\rightarrow$ 0,1977). Le rapport rural/urbain passe de 14,5 à 2,7.

\begin{attention}[Conséquence directe sur le ciblage]
Les variantes ne changent pas seulement le niveau national : elles changent \emph{qui} est
désigné comme pauvre. Une politique calibrée sur « SU3 uniquement » orienterait des moyens vers
des zones urbaines que l'IPM de référence classe parmi les moins pauvres du pays. Le choix de
variante est donc une décision de politique publique, pas un réglage technique.
\end{attention}

\section{Contribution des indicateurs à M0}

La contribution de l'indicateur $j$ s'écrit :

\[
  \mathrm{Contribution}_j = \frac{w_j\,CH_j}{M_0},
\]

où $CH_j$ est le taux de privation censuré : la part de population à la fois pauvre et privée
sur l'indicateur $j$. La somme des contributions vaut 100~\%.

\begin{center}
\small
\begin{tabular}{@{}lrrrr@{}}
\toprule
\textbf{Contribution à $M_0$} & \textbf{Référence} & \textbf{SU3+agri.} &
\textbf{SU3 seul} & \textbf{SU3+NEE} \\
\midrule
Chômage BIT & 1,8~\% & --- & --- & --- \\
SU3 & --- & 9,3~\% & 35,5~\% & 13,0~\% \\
Agriculture de subsistance & 26,8~\% & 24,5~\% & --- & --- \\
NEET approché (NEE) & --- & --- & --- & 19,1~\% \\
\midrule
\textbf{Dimension Emploi} & \textbf{28,6~\%} & \textbf{33,8~\%} & \textbf{35,5~\%} &
\textbf{32,1~\%} \\
Éducation & 40,6~\% & 37,8~\% & 37,1~\% & 39,7~\% \\
Conditions de vie & 29,2~\% & 27,0~\% & 25,5~\% & 26,4~\% \\
Santé & 1,6~\% & 1,4~\% & 1,9~\% & 1,7~\% \\
\bottomrule
\end{tabular}
\end{center}

Deux enseignements. D'abord, l'agriculture de subsistance est, à elle seule, le
\textbf{deuxième contributeur de l'IPM RGPH} après l'alphabétisation --- la retirer ne peut pas
être sans conséquence. Ensuite, la dimension Santé, dotée d'un poids nominal de 25~\%, ne
contribue qu'à 1,6~\%. C'est ce déséquilibre qui amplifie mécaniquement le rôle de l'Emploi
dans le recensement, et qui explique que les écarts entre variantes y soient bien plus marqués
que dans l'enquête.

\section{Tableau de contrôle de l'IPM national de référence}

Le tableau relie règle, fréquence, poids et contribution. Le taux population est non censuré ;
la contribution ne porte que sur les pauvres après censure.

\begin{center}
\scriptsize
\begin{longtable}{@{}p{3.6cm}p{2.6cm}rrr@{}}
\toprule
\textbf{Indicateur} & \textbf{Dimension} & \textbf{Privation pop.} & $w_j$ &
\textbf{Contribution} \\
\midrule
\endfirsthead
\toprule
\textbf{Indicateur} & \textbf{Dimension} & \textbf{Privation pop.} & $w_j$ &
\textbf{Contribution} \\
\midrule
\endhead
Fréquentation scolaire & Éducation & 35,6~\% & 0,0625 & 8,7~\% \\
Année de scolarité & Éducation & 63,2~\% & 0,0625 & 13,4~\% \\
Alphabétisation & Éducation & 70,4~\% & 0,0625 & 13,6~\% \\
État civil & Éducation & 15,2~\% & 0,0625 & 4,9~\% \\
Mortalité dans le ménage & Santé & 1,1~\% & 0,2500 & 1,6~\% \\
Chômage BIT & Emploi & 5,9~\% & 0,1250 & 1,8~\% \\
Agriculture de subsistance & Emploi & 46,0~\% & 0,1250 & 26,8~\% \\
Électricité & Conditions de vie & 11,5~\% & 0,0417 & 2,5~\% \\
Logement & Conditions de vie & 25,0~\% & 0,0417 & 4,7~\% \\
Eau potable & Conditions de vie & 14,3~\% & 0,0417 & 2,8~\% \\
Énergie de cuisson & Conditions de vie & 62,7~\% & 0,0417 & 9,5~\% \\
Toilettes & Conditions de vie & 45,1~\% & 0,0417 & 7,4~\% \\
Biens d'équipement & Conditions de vie & 18,2~\% & 0,0417 & 2,3~\% \\
\midrule
Total des contributions & & & & 100,0~\% \\
\bottomrule
\end{longtable}
\end{center}

Pour mémoire, les deux indicateurs hors IPM national : SU3 prive 26,3~\% de la population et le
NEET approché 30,9~\%.

\section{Forces, limites et risques d'interprétation}

\begin{description}
  \item[SU3 est mieux mesuré sur le RGPH que sur l'EHCVM.] Il est lu dans une variable
  officielle de l'INS plutôt que reconstruit. C'est le principal argument en faveur du
  recensement sur cette dimension précise.

  \item[SU3 est préférable au chômage BIT strict.] Il récupère le découragement et
  l'attachement incomplet au marché du travail : 14,6~\% des 17--40 ans contre 2,6~\%. Il reste
  toutefois limité à cette tranche d'âge.

  \item[L'agriculture mesure une vulnérabilité absente de SU3.] Un chef travaillant dans son
  champ est \emph{occupé} et n'entre pas dans SU3. Sans indicateur de qualité de l'emploi, cette
  réalité --- 46~\% de la population --- disparaît entièrement de la dimension.

  \item[Le proxy agricole du RGPH n'est pas celui de l'EHCVM.] Branche agricole croisée avec
  statut non salarié, contre travail exclusif sur son propre champ. Les deux visent la même
  chose, aucun n'observe la destination de la production.

  \item[SU3 seul donne un poids très fort à une privation unique.] Chaque ménage SU3 reçoit 0,25
  d'un coup, tout en ignorant la qualité des emplois existants. Le report massif vers la
  vulnérabilité (32,6~\%) en est le symptôme.

  \item[Le NEE n'est pas un NEET complet.] La formation en cours manque, le taux est donc une
  borne supérieure. L'indicateur ne concerne par ailleurs que les 54,3~\% de ménages comprenant
  un jeune de 15 à 24 ans.

  \item[La dimension Santé est structurellement faible.] Un poids de 0,25 pour une contribution
  de 1,6~\%. Toute conclusion tirée de l'IPM RGPH doit tenir compte de ce que la mortalité sur
  12 mois ne peut pas porter, à elle seule, un quart de l'indice.
\end{description}

\section{Recommandation}

\begin{enumerate}
  \item Retenir comme variante principale \textbf{SU3 + dépendance agricole non salariée du
  chef}, avec deux indicateurs de poids 0,125. C'est la seule qui couvre à la fois la
  sous-utilisation de la main-d'{\oe}uvre et la précarité de l'emploi existant.
  \item Qualifier explicitement l'indicateur agricole de \textbf{proxy}, car la destination de la
  production n'est pas observée.
  \item Conserver \textbf{SU3 uniquement} comme test de sensibilité, en documentant que la baisse
  de $M_0$ est un effet de suppression d'indicateur, pas une amélioration.
  \item Conserver \textbf{SU3 + NEE} comme scénario exploratoire, sans l'appeler NEET dans une
  publication officielle tant que la formation en cours n'est pas mesurée.
  \item Ne jamais présenter ces quatre chiffres côte à côte sans le tableau par milieu : c'est
  là que se joue la conséquence réelle du choix.
  \item Au prochain recensement, ajouter une question sur la destination principale de la
  production agricole et une question sur la formation suivie actuellement.
\end{enumerate}

\section{Contrôles réalisés pour sécuriser les chiffres}

\begin{itemize}
  \item auto-contrôles unitaires des étapes 01, 02 et 05, étendus aux nouvelles colonnes : un
  inactif n'entre ni dans SU1 ni dans SU3 ; une personne de la main-d'{\oe}uvre potentielle
  entre dans SU3 mais pas dans SU1 ; un jeune scolarisé n'est pas NEET ; un âge « ne sait pas »
  n'entre dans aucune tranche ;
  \item contrôle que tout ménage privé au chômage BIT est également privé selon SU3 : 0 exception
  sur 196\,019 ménages ;
  \item \col{assert} de l'étape 05 vérifiant que la variante nationale reproduit à l'identique la
  matrice censurée produite indépendamment par l'étape 04 ;
  \item contrôle que la somme des lignes censurées redonne exactement $c_i(k)$ : écart maximal
  nul ;
  \item contrôle de décomposabilité : pour chaque variable de désagrégation, les $M_0$ de groupe
  pondérés par la part de population redonnent le $M_0$ national ;
  \item vérification que la somme des poids vaut 1 à $10^{-9}$ près dans chacune des six
  variantes ;
  \item vérification que l'IPM national du RGPH est inchangé après l'ajout de SU3 et du NEET :
  $H = 38{,}8~\%$, $A = 0{,}4287$, $M_0 = 0{,}1664$, identiques aux valeurs publiées ;
  \item distribution complète de \col{Statut\_OQPtbb} relue sur les 27\,984\,405 individus du
  fichier pour confirmer le codage des quatre modalités.
\end{itemize}

\begin{retenir}[Conclusion]
Les variantes Emploi sont désormais calculables sur le recensement, et SU3 y est même mieux
mesuré que dans l'enquête puisqu'il est lu dans une variable officielle de l'INS. Le choix
restant n'est pas informatique : il consiste à décider quelle conception de la privation
d'emploi doit entrer dans l'IPM. Sur le RGPH, l'enjeu est plus lourd que sur l'EHCVM, parce que
la dimension Santé y est presque muette et que l'Emploi y pèse d'autant plus. Les résultats
conduisent à préférer SU3 accompagné de l'indicateur agricole, sous un libellé prudent.
\end{retenir}

\section*{Fichiers de résultats associés}

\begin{itemize}
  \item \col{indices\_ipm\_su3\_rgph2021.xlsx} : SU3 + agriculture ;
  \item \col{indices\_ipm\_su3\_seul\_rgph2021.xlsx} : SU3 uniquement ;
  \item \col{indices\_ipm\_su3\_neet\_rgph2021.xlsx} : SU3 + NEE ;
  \item \col{contributions\_ipm\_su3*\_rgph2021.csv} : contributions par indicateur ;
  \item chaque classeur contient les résultats nationaux, les contributions, les poids et les
  désagrégations par milieu, région, département, sous-préfecture et taille du ménage.
\end{itemize}

\end{document}
